\documentclass[10pt,a4paper]{article}

\usepackage[utf8]{inputenc}
\usepackage[T1]{fontenc}
\usepackage{lmodern}
\usepackage[slovak]{babel}
\usepackage[a4paper,margin=1.8cm]{geometry}
\usepackage{booktabs}
\usepackage{tabularx}
\usepackage{array}
\usepackage{longtable}
\usepackage{multirow}
\usepackage{enumitem}
\usepackage[hidelinks]{hyperref}
\usepackage[protrusion=true,expansion=false]{microtype}
\usepackage{parskip}

% Spoločné údaje projektu — pri zmene upravte iba tento súbor.
% Dôležité: schválené sú slovenský názov a anotácia. EN názov/akronym sú návrhy;
% dátumy, rozpočet a PI sú pracovné predpoklady semestrálnej verzie, nie údaje z APVV výzvy.
\newcommand{\projectTitleSK}{Detekcia falošných online recenzií na základe jazykových a behaviorálnych príznakov}
\newcommand{\projectTitleEN}{Detection of Fake Online Reviews Based on Linguistic and Behavioral Features}
\newcommand{\projectAcronym}{FAKERADAR}
\newcommand{\projectId}{Zatiaľ nepridelené (študentský návrh)}
\newcommand{\projectDate}{01.11.2026}
\newcommand{\projectStart}{01/2027}
\newcommand{\projectEnd}{12/2029}

\newcommand{\applicant}{Slovenská technická univerzita v Bratislave, Fakulta informatiky a informačných technológií}
\newcommand{\applicantShort}{FIIT STU}
\newcommand{\applicantAddress}{Ilkovičova 6276/2, 842 16 Bratislava 4}
\newcommand{\applicantICO}{(doplní sa)}
\newcommand{\applicantDIC}{(doplní sa)}
\newcommand{\applicantICDPH}{(doplní sa)}
\newcommand{\applicantEmail}{info@fiit.stuba.sk}
\newcommand{\applicantPhone}{+421 2 210 44 111}
\newcommand{\statutory}{prof. Ing. Ivan Kotuliak, PhD.}
\newcommand{\statutoryRole}{dekan}
\newcommand{\piName}{Ing. Ivan Kapustík}
\newcommand{\piEmail}{ivan.kapustik@stuba.sk} % e-mail nie je overený; pred odovzdaním skontrolovať

\newcommand{\budgetAPVV}{86\,160}
\newcommand{\budgetCoFin}{0}
\newcommand{\budgetTotal}{86\,160}

% Anotácia = presné znenie schváleného zamerania (odovzdané 04.10.2026, hodnotené cvičiacim 5/5 b);
% zdroj: old_zameranie_clanok.tex (abstrakt).
\newcommand{\projectAnnotation}{%
Online recenzie výrazne ovplyvňujú rozhodovanie spotrebiteľov, no ich dôveryhodnosť narúšajú falošné recenzie generované firmami, platenými recenzentmi či automatizovanými nástrojmi. Tento článok prináša prehľad súčasných metód detekcie falošných recenzií so zameraním na porovnanie jazykových (textových) a behaviorálnych (na správaní recenzenta založených) príznakov.}


\setlength{\parindent}{0pt}
\setlength{\emergencystretch}{3em}
\newcommand{\formsection}[1]{%
  \section*{\raggedright #1}%
  \vspace{-0.4em}%
}
\newcommand{\field}[2]{\textbf{#1} & #2 \\}

\begin{document}

\begin{titlepage}
\centering
\vspace*{1.5cm}
{\Large\bfseries Žiadosť o riešenie projektu aplikovaného výskumu a vývoja}\\[0.4em]
{\large Formuláre VV-A až VV-E}\\[2.5cm]
{\LARGE\bfseries \projectTitleSK}\\[0.8em]
{\large \projectTitleEN}\\[2cm]
\begin{tabular}{rl}
\textbf{Akronym:} & \projectAcronym\\
\textbf{Evidenčné číslo:} & \projectId\\
\textbf{Dátum podania:} & \projectDate\\
\textbf{Žiadateľ:} & \parbox[t]{9.2cm}{\raggedright \applicant}\\
\end{tabular}
\vfill
{\footnotesize
Dokument je modelovým návrhom vypracovaným v rámci semestrálneho projektu v predmete Metódy inžinierskej práce. Ekonomické údaje a niektoré personálne údaje sú modelové; zástupné hodnoty označené ako „(doplní sa)“ sa doplnia vo finálnej verzii.
}
\end{titlepage}

\newpage

% ---------------------------------------------------------------
\formsection{VV-A1 — Základné informácie o projekte}
\begin{tabularx}{\textwidth}{@{}>{\raggedright\arraybackslash}p{5.7cm} X@{}}
\toprule
\textbf{Pole} & \textbf{Údaj} \\
\midrule
01 Evidenčné číslo projektu & \projectId \\
02 Dátum podania & \projectDate \\
03 Názov projektu (SK) & \projectTitleSK \\
03 Názov projektu (EN) & \projectTitleEN \\
04 Akronym projektu & \projectAcronym \\
05 Odbor vedy a techniky & Informatika — počítačové vedy a informačné systémy \\
06 Charakter výskumu & Aplikovaný výskum \\
07 Začiatok riešenia projektu & \projectStart \\
08 Koniec riešenia projektu & \projectEnd \\
09 Anotácia (max.\ 2\,000 znakov) & \projectAnnotation \\
10 Žiadateľská organizácia & \applicant \\
11 Požadované finančné prostriedky z APVV & \budgetAPVV\ EUR \\
12 Spolufinancovanie projektu & \budgetCoFin\ EUR \\
13 Celkové náklady na projekt & \budgetTotal\ EUR \\
\bottomrule
\end{tabularx}

\smallskip
\emph{Poznámka k pôvodu údajov: slovenský názov a anotácia sú prevzaté zo schváleného zamerania. Anglický názov a akronym sú návrhy, keďže v odovzdanom zameraní chýbali. Termín 01/2027--12/2029, TRL a finančné sumy sú plánovacie predpoklady tímu pre 36-mesačný modelový projekt.}

\newpage
% ---------------------------------------------------------------
\formsection{VV-A2 — Základné informácie o riešiteľských organizáciách}

\subsection*{Žiadateľ}
\begin{tabularx}{\textwidth}{@{}>{\raggedright\arraybackslash}p{5.7cm} X@{}}
\toprule
\textbf{Pole} & \textbf{Údaj} \\
\midrule
01 Názov organizácie & \applicant \\
02 Skrátený názov & \applicantShort \\
03 Adresa organizácie & \applicantAddress \\
04 IČO & \applicantICO \\
05 Príslušnosť k rezortu & Školstvo, veda, výskum a šport \\
06 Forma hospodárenia & Verejná vysoká škola \\
07 Kontaktná osoba & Matvii Nosenko, \texttt{xnosenko@stuba.sk}, \applicantPhone \\
08 Štatutárny zástupca I & \statutory, \statutoryRole \\
09 Štatutárny zástupca II & — \\
\bottomrule
\end{tabularx}

\subsection*{Spoluriešiteľská organizácia}
Bez spoluriešiteľskej organizácie. Projekt je riešený na jednej organizácii; pilotní partneri z praxe budú určení počas riešenia a ich účasť sa upraví zmluvne.

\newpage
% ---------------------------------------------------------------
\formsection{VV-A3 — Zoznam riešiteľov}

\subsection*{01 Zamestnanci priamo sa podieľajúci na riešení projektu}
{\footnotesize
\setlength{\tabcolsep}{3.5pt}
\begin{longtable}{@{}p{0.5cm} p{2.9cm} p{1.1cm} p{2.5cm} p{2.2cm} p{1.7cm} p{2.2cm} p{1.2cm}@{}}
\toprule
\# & Meno a priezvisko & Tituly & Pracovné zaradenie & Dátum narodenia & IČO organizácie & Počet hodín/rok & Vek $\leq$35? \\
\midrule
\endhead
1 & Ivan Kapustík & Ing. & zodpovedný riešiteľ & (doplní sa) & \applicantICO & 600 (modelový údaj) & overiť \\
2 & Matvii Nosenko & Bc. & študent 1. stupňa -- riešiteľ & (doplní sa) & \applicantICO & 100 & overiť \\
3 & Volodymyr Nedoluha & Bc. & študent 1. stupňa -- riešiteľ & (doplní sa) & \applicantICO & 100 & overiť \\
4 & Danylo Nikolenko & Bc. & študent 1. stupňa -- riešiteľ & (doplní sa) & \applicantICO & 100 & overiť \\
5 & Rostyslav Mykytiuk & Bc. & študent 1. stupňa -- riešiteľ & (doplní sa) & \applicantICO & 100 & overiť \\
\bottomrule
\end{longtable}
}

\vspace{0.5em}
Poznámka: dátum narodenia je v žiadosti povinný; v predbežnej verzii je zástupný. Členovia tímu sú študenti prvého stupňa VŠ (Bc.); orientačný limit je 100 hodín ročne na študenta. Dátum narodenia je povinný a doplní sa; vekový status do 35 rokov sa určí podľa neho. Kapacita PI 600 hodín a podporná kapacita 300 hodín sú pracovné odhady. Štyria študenti tvoria 80\,\% z piatich uvedených riešiteľov. Hodiny v tejto tabuľke sú pracovné predpoklady pre predbežný plán, nie potvrdený pracovný výkaz ani mzdový výpočet.

\subsection*{02 Ostatní zamestnanci}
Modelovo predpokladáme administratívnu a technickú podporu organizácie v rozsahu približne 300 hodín ročne (modelový údaj).

\subsection*{03 Spolu}
Počet priamo uvedených riešiteľov: 5. Ich pracovná kapacita je v tomto modeli 1\,000 hodín ročne (PI 600 + 4 študenti × 100). Po pripočítaní modelovej podpornej kapacity 300 hodín je plánovací súčet približne 1\,300 hodín ročne.

\newpage
% ---------------------------------------------------------------
\formsection{VV-A4 — Základné informácie o zodpovednom riešiteľovi}
\begin{tabularx}{\textwidth}{@{}>{\raggedright\arraybackslash}p{5.7cm} X@{}}
\toprule
\textbf{Pole} & \textbf{Údaj} \\
\midrule
01 Meno a priezvisko & \piName \\
02 Pohlavie & muž \\
03 Telefón & \applicantPhone\ (ústredňa FIIT STU; osobný kontakt sa doplní vo finálnej verzii) \\
04 E-mail & \texttt{\piEmail} \\
05 Zodpovedný riešiteľ je mladým vedeckým pracovníkom & nie \\
06 ID výskumníka (vedecká databáza) & nevypĺňa sa (nie je povinné) \\
07 Prehľad projektov zodpovedného riešiteľa za posledných 5 rokov & doplní sa; uvádzajú sa ľubovoľné projekty, najlepšie tie, ktoré riešiteľ viedol (štruktúra: názov, schéma, roky realizácie, rozpočet, pozícia) \\
\bottomrule
\end{tabularx}

\subsection*{Čestné vyhlásenie zodpovedného riešiteľa}
Ja, dolu podpísaný/á \piName, čestne vyhlasujem, že údaje uvedené v častiach „Základné informácie o zodpovednom riešiteľovi“ a v štruktúrovanom profesijnom životopise sú pravdivé.

Ja, dolu podpísaný/á \piName, čestne vyhlasujem, že v rámci riešenia projektu budú zo strany žiadateľa rešpektované všetky zásady etických princípov vyplývajúce z európskej a národnej legislatívy.

\vspace{0.8em}
Podpis zodpovedného riešiteľa: \dotfill\ Miesto: Bratislava, dátum: \projectDate

\newpage
% ---------------------------------------------------------------
\formsection{VV-B — Ciele, zámery a výstupy projektu}

\begin{tabularx}{\textwidth}{@{}>{\raggedright\arraybackslash}p{5.7cm} X@{}}
\toprule
\textbf{Pole} & \textbf{Údaj} \\
\midrule
01 Kľúčové slová & falošné recenzie; recenzný spam; spracovanie prirodzeného jazyka; behaviorálna analýza; hybridný model; vysvetliteľnosť; slovenský jazyk \\
02 Ciele projektu & viď zoznam nižšie \\
03 Zaradenie projektu do TRL & pracovný návrh: TRL 3 na začiatku, TRL 6 na konci \\
03 Zdôvodnenie zaradenia & Návrhový predpoklad: začiatok pri overovaní konceptu a koniec funkčným prototypom demonštrovaným v relevantnom prostredí; tím musí úroveň potvrdiť podľa definícií TRL \\
03 Zaradenie výstupov do TRL & pracovný návrh: dataset a metodika TRL 4; model TRL 5; prototyp API TRL 6 \\
04 Využitie výsledkov — odberateľ (realizátor) je žiadateľ & nasadenie detektora vo vlastných výskumných, vzdelávacích a konzultačných aktivitách FIIT STU \\
04 Využitie výsledkov — iný odberateľ & prevádzkovatelia e-shopov a recenzných platforiem; orgány ochrany spotrebiteľa a regulátor \\
05 Prehľad výstupov & viď tabuľku nižšie \\
\bottomrule
\end{tabularx}

\subsection*{Ciele projektu}
\begin{enumerate}[label=O\arabic*., leftmargin=*]
  \item Vytvoriť otvorený slovensko-anglický anotovaný dataset s pracovným cieľom 15\,000 recenzií a mierou zhody anotátorov $\kappa \geq 0{,}7$.
  \item Vyvinúť a vyhodnotiť jazykový model dosahujúci F1 najmenej 0{,}80 v rámci domény.
  \item Vyvinúť behaviorálny a grafový model dosahujúci F1 najmenej 0{,}75 a odhaľujúci koordinované skupiny.
  \item Skonštruovať hybridný model s F1 najmenej 0{,}85 v rámci domény a 0{,}75 pri prenose na novú doménu.
  \item Dodať prototyp API s vysvetliteľnými výstupmi a vyhodnotiť ho v partnerskom pilote, ak sa zabezpečí prístup; inak vykonať offline test.
\end{enumerate}

\subsection*{Prehľad plánovaných výstupov}
\begin{tabularx}{\textwidth}{@{}p{1.2cm} X p{3.0cm}@{}}
\toprule
\# & Výstup & Termín \\
\midrule
V1 & Metodika anotácie a anotačné guidelines & M6 \\
V2 & Otvorený dataset SK/EN & M12 \\
V3 & Jazykový model a vyhodnotenie & M18 \\
V4 & Behaviorálny a grafový model & M24 \\
V5 & Hybridný model s vysvetliteľnosťou & M30 \\
V6 & Prototyp API a pilotná správa & M36 \\
V7 & Dve odborné publikácie a jedna správa pre prax & M36 \\
\bottomrule
\end{tabularx}

\smallskip
\emph{Ciele, prahy F1, cieľová veľkosť datasetu a TRL sú predbežné návrhové hodnoty tímu; nie sú predpísané webom predmetu. Ich dôvody a plán overenia sú opísané vo VV-F, časti 1.4.}

\newpage
% ---------------------------------------------------------------
\formsection{VV-C — Rozpočet projektu v EUR}

\begin{table}[h!]
\centering
\small
\begin{tabularx}{\textwidth}{@{}X r r r r@{}}
\toprule
\textbf{Kategória} & \textbf{2027} & \textbf{2028} & \textbf{2029} & \textbf{Spolu} \\
\midrule
01 Bežné priame náklady & & & & \\
02 Mzdové náklady a ostatné osobné náklady & 17\,200 & 17\,200 & 17\,200 & 51\,600 \\
03 Zdravotné a sociálne poistenie & 6\,020 & 6\,020 & 6\,020 & 18\,060 \\
04 Cestovné náklady & 1\,000 & 1\,000 & 1\,000 & 3\,000 \\
05 Materiál & 1\,000 & 1\,000 & 1\,000 & 3\,000 \\
06 Služby & 2\,000 & 2\,000 & 2\,000 & 6\,000 \\
07 Energie, vodné, stočné, komunikácie & 500 & 500 & 500 & 1\,500 \\
08 Náklady na popularizáciu & 1\,000 & 1\,000 & 1\,000 & 3\,000 \\
09 Bežné nepriame náklady & 0 & 0 & 0 & 0 \\
\midrule
10 Bežné náklady spolu & 28\,720 & 28\,720 & 28\,720 & 86\,160 \\
\midrule
Celkové náklady z APVV & 28\,720 & 28\,720 & 28\,720 & 86\,160 \\
Spolufinancovanie & 0 & 0 & 0 & 0 \\
\midrule
Celkové náklady & 28\,720 & 28\,720 & 28\,720 & 86\,160 \\
\bottomrule
\end{tabularx}
\end{table}

\subsection*{Predpoklady modelu}
Mzdy sú vypočítané z predpokladaných hodín: PI 600 h/rok × 22 EUR/h hrubo a štyria študenti 4 × 100 h/rok × 10 EUR/h hrubo, spolu 17\,200 EUR ročne. Odvody sú 35 % miezd. Ostatné položky sú odhady bez potvrdených ponúk. Projekt je študentský, preto sa neplánuje spolufinancovanie.

Ide o modelový rozpočet študentského projektu v predmete MIP, nie o rozpočet reálneho grantu ani o schválený rozpočet FIIT. Rozdelenie medzi roky je rovnomerné: 28\,720 EUR ročne, spolu 86\,160 EUR (3 × 28\,720 = 86\,160). Odvody sú modelované ako 35\,\% mzdových nákladov (6\,020 EUR ročne); pred finálnou verziou treba overiť hodiny a sadzby.

\newpage
% ---------------------------------------------------------------
\formsection{VV-D — Harmonogram projektu}

\begin{tabularx}{\textwidth}{@{}p{0.8cm} p{1.8cm} p{1.8cm} X p{1.8cm}@{}}
\toprule
\# & Začiatok & Koniec & Názov etapy & Míľnik \\
\midrule
1 & 01/2027 & 06/2027 & WP1 Manažment projektu a príprava & M6 \\
2 & 01/2027 & 12/2027 & WP2 Zber, anotácia a sprístupnenie dát & M12 \\
3 & 07/2027 & 06/2028 & WP3 Jazykové modely & M18 \\
4 & 01/2028 & 12/2028 & WP4 Behaviorálne a grafové modely & M24 \\
5 & 07/2028 & 06/2029 & WP5 Hybridný model a vysvetliteľnosť & M30 \\
6 & 01/2029 & 12/2029 & WP6 Pilotné overenie a evaluácia & M36 \\
7 & 01/2027 & 12/2029 & WP7 Diseminácia a komunikácia & priebežne \\
\bottomrule
\end{tabularx}

\subsection*{Míľniky}
\begin{itemize}[leftmargin=*]
  \item M6: schválená metodika anotácie a etický rámec.
  \item M12: dataset SK/EN s minimálne 15\,000 recenziami.
  \item M18: jazykový model a prvá publikácia.
  \item M24: behaviorálny a grafový model.
  \item M30: hybridný model a prototyp API.
  \item M36: pilotná správa, finálne vyhodnotenie a druhá publikácia.
\end{itemize}

\subsection*{Zoznam prístrojov v projekte}
\begin{tabularx}{\textwidth}{@{}X X@{}}
\toprule
\textbf{Názov} & \textbf{Popis} \\
\midrule
Výpočtový server s GPU (FIIT STU) & trénovanie a evaluácia transformerových modelov; zdieľaná infraštruktúra fakulty \\
Pracovné stanice riešiteľov & predspracovanie dát, anotácia, exploratívna analýza \\
Úložisko a verzovací server & uchovanie datasetu, kódu a experimentov (Git, zálohy) \\
\bottomrule
\end{tabularx}

\newpage
% ---------------------------------------------------------------
\formsection{VV-E — Čestné vyhlásenie štatutárneho zástupcu žiadateľskej organizácie}

Ja, dolu podpísaný/á \statutory, \statutoryRole{} organizácie \applicant, záväzne vyhlasujem, že:

\begin{itemize}[leftmargin=*]
  \item organizácia má platné osvedčenie o spôsobilosti vykonávať výskum a vývoj v zmysle podmienok stanovených vo verejnej výzve v súlade s \S~26a zákona č.~172/2005 Z.~z. v znení neskorších predpisov a je evidovaná v zozname spôsobilých osôb vykonávať výskum a vývoj;
  \item všetky informácie obsiahnuté v dokumentácii návrhu sú pravdivé;
  \item organizácia má, resp. zabezpečí vlastné zdroje na spolufinancovanie projektu;
  \item predložený návrh projektu nebol a nie je financovaný z iných zdrojov (národných alebo zahraničných);
  \item organizácia neposkytne žiadnu nepriamu štátnu pomoc v zmysle špecifických podmienok výzvy;
  \item organizácia nie je daňovým dlžníkom;
  \item organizácia nie je dlžníkom poistného na sociálnom poistení (dôchodkovom, nemocenskom, garančnom a úrazovom poistení, poistení v nezamestnanosti, poistení do rezervného fondu solidarity) a dlžníkom príspevkov na starobné dôchodkové sporenie;
  \item organizácia nie je dlžníkom poistného na zdravotnom poistení (za každú zdravotnú poisťovňu);
  \item organizácia nie je v likvidácii;
  \item voči organizácii nie je vedený výkon rozhodnutia;
  \item voči organizácii nebolo začaté konkurzné/reštrukturalizačné konanie;
  \item na majetok organizácie nie je vyhlásený konkurz;
  \item voči organizácii nebol zamietnutý návrh na vyhlásenie konkurzu pre nedostatok majetku;
  \item organizácia neporušila zákaz nelegálnej práce a nelegálneho zamestnávania podľa osobitého predpisu;
  \item organizácia má vysporiadané finančné vzťahy so štátnym rozpočtom;
  \item organizácia bude súčinná pri poskytovaní dokladov, ktoré preukazujú splnenie podmienok na poskytnutie dotácie v zmysle zákona č.~523/2004 Z.~z.;
  \item som si vedomý/á povinnosti poskytnutia informácií o výskume a vývoji, na ktorý boli poskytnuté finančné prostriedky zo štátneho rozpočtu, za účelom zverejnenia v súlade s platnou legislatívou SR;
  \item som si vedomý/á povinnosti zverejnenia informácií o technickej infraštruktúre výskumu a vývoja na centrálnom informačnom portáli;
  \item som si vedomý/á skutočnosti, že APVV zverejní na webovej stránke \url{www.apvv.sk} rozhodnutie o žiadostiach odporučených aj neodporučených orgánom APVV na financovanie;
  \item som si vedomý/á povinnosti popularizovať výsledky riešenia projektu s cieľom zrozumiteľným spôsobom informovať verejnosť o prínosoch výsledkov a výstupov výskumu a vývoja;
  \item predložený návrh projektu po svojej obsahovej stránke nebude financovaný z národných zdrojov iných grantových schém ako z príslušnej grantovej schémy APVV.
\end{itemize}

Ako štatutárny zástupca svojím podpisom zodpovedám za správnosť a pravdivosť údajov uvedených v tejto elektronickej žiadosti. V prípade vyzvania zo strany APVV sa zaväzujem nahradiť toto čestné vyhlásenie aktuálnymi potvrdeniami príslušných úradov. Som si vedomý/á, že v prípade zistenia, že údaje uvedené v predmetnom vyhlásení nie sú pravdivé, žiadosť bude vyradená alebo zmluva o poskytnutí prostriedkov nebude podpísaná, prípadne dôjde k odstúpeniu od zmluvy.

\vspace{1.2em}
{\footnotesize
\begin{tabularx}{\textwidth}{@{}p{7.2cm} X@{}}
\toprule
\textbf{Názov žiadateľskej organizácie} & \applicant \\
\textbf{Meno štatutárneho zástupcu (I)} & \statutory \\
\textbf{V zastúpení (uveďte čitateľne meno)} & --- \\
\textbf{Podpis štatutárneho zástupcu} & \dotfill \\
\textbf{Miesto} & Bratislava \\
\textbf{Dátum} & \projectDate \\
\bottomrule
\end{tabularx}
}

\end{document}
